% Generated draft: requires mathematical and source review.
\documentclass[11pt]{article}
\usepackage[margin=25mm]{geometry}
\usepackage{fontspec}
\setmainfont{Latin Modern Roman}
\usepackage{amsmath,amssymb,mathtools,unicode-math}
\setmathfont{Latin Modern Math}
\usepackage{enumitem,needspace,xurl,hyperref}
\setlength{\parindent}{0pt}
\setlength{\parskip}{5pt}
\setlength{\emergencystretch}{3em}
\newcommand{\N}{\mathbb{N}}
\newcommand{\Z}{\mathbb{Z}}
\newcommand{\Q}{\mathbb{Q}}
\newcommand{\R}{\mathbb{R}}
\newcommand{\C}{\mathbb{C}}
\newcommand{\F}{\mathbb{F}}
\newcommand{\A}{\mathbb{A}}
\newcommand{\PP}{\mathbb{P}}
\newcommand{\E}{\mathbb{E}}
\newcommand{\eps}{\varepsilon}
\newcommand{\dd}{\,\mathrm{d}}
\begin{document}
\section*{UnsolvedMath ID: 20001925}
\textbf{Source problem code:} catalogue entry
\subsection{Definitions and mathematical statement}
Let $M$ be a smooth closed manifold of dimension $n\ge2$, and let $g$ be a smooth semi-Riemannian metric of indefinite signature $(p,q)$ with $p,q\ge1$. A null geodesic is a nonconstant affinely parametrized geodesic $\gamma$ with $g(\dot\gamma,\dot\gamma)=0$. Null completeness means that every such geodesic is defined for all $t\in\R$; geodesic completeness, assumed in the original question, is stronger.

\textbf{Literal Question 9.2.1.} Can such a complete $(M,g)$ admit a nonzero smooth one-form $\eta$ for which
\[
 \forall\gamma\ \exists a_\gamma\in\R\setminus\{0\}\ 
 \exists b_\gamma\in\R\ \forall t\in\R,
 \qquad \eta_{\gamma(t)}(\dot\gamma(t))=a_\gamma t+b_\gamma,
\]
where $\gamma$ ranges over all nonconstant affinely parametrized null geodesics? The constants may depend on the geodesic. ``Grows linearly in both directions'' is understood by this displayed affine formula, not as positivity in both time directions. The source-verification field identifies the OCR string $C_1\,6=0$ as $C_1\ne0$.

\textbf{Supplied negative-answer claim.} The aid asserts that this property is impossible already under null completeness and for $\eta$ of class $C^1$. To state its compactness reduction, choose any auxiliary Riemannian metric $h$ and put
\[
 \mathcal N_h=\{(x,v)\in TM:g_x(v,v)=0,\ |v|_h=1\},
 \qquad A_\eta(x,v)=(\nabla_v\eta)(v).
\]
The claimed theorem says that, if evaluation of $\eta$ is affine on every complete null geodesic, then $A_\eta$ vanishes at some point of $\mathcal N_h$. Equivalently, at least one such geodesic has zero slope. The aid describes the obstruction using compactness of $\mathcal N_h$, a uniform positive lower bound for $|A_\eta|$ if it has no zeros, and conservation of $A_\eta(\gamma,\dot\gamma)$ along each affine-evaluation orbit. This is a supplied candidate theorem, not a verified literature citation.

\textbf{Projective-rigidity context and quantifier limitation.} The source states the semi-Riemannian projective Lichnerowicz--Obata conjecture: if a connected Lie group acts on a complete $(M^n,g)$, $n\ge2$, by diffeomorphisms preserving unparametrized geodesics, then its action is isometric or $g$ has constant sectional curvature. The source reports the Riemannian case in [Matveev2005, Matveev2007] and the closed two-dimensional case in [Matveev2012b], while noting difficulties in indefinite dimension at least three.

For a projective vector field, identified with its metric-dual one-form $\eta$, the source quotes a second-order projective-field equation involving an auxiliary one-form $\lambda$. Its null contraction is asserted to imply
\[
 \frac{\dd^2}{\dd t^2}\eta(\dot\gamma(t))=0.
\]
The printed extraction of the full indexed equation is damaged; its precise derivative and trace conventions are not reconstructed here. The subsequent conformal-field argument requires zero slope on \emph{every} null geodesic, which would give
\[
 \nabla_i\eta_j+\nabla_j\eta_i=\rho g_{ij}
\]
for some function $\rho$. The negative answer to the literal existence question supplies only \emph{some} zero-slope null direction for a given affine-evaluation form. Therefore the source's implication to full projective rigidity does not follow from that negative answer alone. Its text also says ``positive'' at one point where its surrounding discussion says ``negative''; both inconsistencies are retained as qualifications. The final step invokes Weyl's assertion that a transformation both projective and conformal is a homothety; on a closed manifold a global positive homothety factor is one by volume preservation.

The source additionally asks whether the velocity of almost every geodesic stays in a bounded subset of $TM$ for complete semi-Riemannian metrics on closed manifolds. On compact $M$, boundedness can be specified as $\sup_{t\in\R}|\dot\gamma(t)|_h<\infty$. A measure on initial conditions is not specified, and an ambient almost-everywhere assertion cannot simply be substituted for a null-cone assertion. The source presents this as another possible route to its conjecture, not as an established theorem.

The extract ends with the start of Section 9.3, communicated by H.~Baum: closed Riemannian manifolds are geodesically complete, whereas closed manifolds with indefinite metrics can be incomplete. It supplies no further complete Section 9.3 problem.
\subsection{Short English statement}
The literal question demands a nonzero affine slope on every complete null geodesic. The supplied aid claims that compactness rules this out even with only null completeness. This gives a zero-slope direction, which is insufficient for the source's proposed deduction that all projective vector fields are conformal.
\subsection{Sources}
\begin{itemize}
\item Local UnsolvedMath registry, permanent ID 20001925, source code catalogue entry; canonical \texttt{aim-geometry-notes.json}, notes[262].
\item Keith Burns and Vladimir Matveev, \emph{Open problems and questions about geodesics}, Question 9.2.1, AIM source \url{https://aimath.org/pastworkshops/geodesicsproblems.pdf}. The record reports verification at printed p.~22, PDF page 21, including $C_1\ne0$; this is source-reported verification.
\item Supplied labels [Matveev2005], [Matveev2007], [Matveev2012b], and [Weyl1921]; no full bibliographic entries are supplied. Section 9.3 is attributed to H.~Baum.
\item Research attempt 1, registry update 2026-08-14: full solution, problem status at run ``refuted'', rigor ``proved'', novelty confidence low. The aid is explicitly machine-generated, not peer reviewed, and must be independently verified before citation.
\end{itemize}
\end{document}
